\begin{definition}[The order on symmetric matrices, (2.1)]
Let $m\in\mathbb N$ and $A,B\in\mathbb R^{m\times m}$. We write
$\mathrm{MatLe}\ A\ B$, read $A\le B$, if
\[
 \langle A\xi,\xi\rangle\ \le\ \langle B\xi,\xi\rangle\qquad\text{for all }\xi\in\mathbb R^m,
 \qquad\text{where}\quad
 \langle A\xi,\xi\rangle=\sum_{i=1}^m\sum_{j=1}^m A_{i,j}\,\xi_j\,\xi_i .
\]
This is the order (2.1) of the paper's notation section. No symmetry is required of
$A$ or $B$ in the definition; the relation only sees their symmetric parts, and it is
applied in the paper exclusively to symmetric matrices, where it is the usual partial
order. Two uses fix the strength needed: the trace monotonicity (2.2)
($\mathrm{Tr}(AY)\le\mathrm{Tr}(AX)$ whenever $A\ge0$ and $Y\le X$), and the comparison
$\mathcal X\mathcal X^{s}u(x_0)\le\mathcal X\mathcal X^{s}v(x_0)$ of the symmetrized
horizontal second derivatives of a semi-convex $u$ and a semi-concave $v$ at a local
maximum point $x_0$ of $u-v$, which opens the proof of Proposition Proposition~3.3.
\end{definition}
